\section{Limitations and Ethical Considerations}
\label{sec:limits-ethics}

\textbf{Limitations.} We acknowledge several limitations. (1)~\emph{Agent harness.} We use a minimal ReAct harness as the primary configuration; our harness ablation (\Cref{sec:res-harness}) on both synthetic and UCI tasks shows that silent failures persist under single-shot and few-shot variants. However, the harness significantly affects the failure distribution, so we report harness as a benchmark dimension. Production agents (e.g., code-interpreter, AutoGen) remain untested. (2)~\emph{Model scope.} DeepSeek-R1 and Qwen3-max show 0\% success on UCI tasks; our enriched manifest (\texttt{manifest\_v3.csv}) reveals these are primarily infrastructure failures (API/parse/format issues), not genuine model inability. We report \texttt{SR$^v$} (success rate among valid runs) separately from raw SR to address this. (3)~\emph{Taxonomy is under-populated; the current benchmark is effectively binary.} This is our most important limitation. Although we define a five-stage taxonomy (Analytical-Plan, Code-Generation, Runtime, Output-Mismatch, Answer-Error), the labeled data is dominated by two classes: \texttt{output\_mismatch} (2{,}202) and \texttt{runtime} (600), with only 1 \texttt{analytical\_plan} case and 0 \texttt{code\_generation} cases. The five-stage taxonomy should therefore be read as a \emph{framework design} rather than an empirically validated five-way benchmark: in its current form AgentFail is a binary runtime-vs-output\_mismatch (loud-vs-silent) diagnostic benchmark. The absence of Code-Generation labels stems from the rule-based classifier, which cannot reliably separate code-level operation errors from downstream answer-extraction errors without ground-truth reference-path annotations, and thus merges them into \texttt{output\_mismatch}. Populating Analytical-Plan, Code-Generation, and Answer-Error with 100--200 gold-annotated traces each, and reporting per-class confusion matrices and macro-F1, is required future work before the five-stage claim can be made. (3b)~\emph{Label reliability.} Three human annotators achieved Fleiss' $\kappa = 0.860$ (strong agreement) on 47 traces, but this measures human--human agreement, not classifier correctness. The rule-based classifier achieves only 71.8\% accuracy against adjudicated gold labels ($\kappa = 0.438$), with a systematic bias toward \texttt{output\_mismatch} over \texttt{runtime}; the roughly 3{,}940 automatically labeled failures the benchmark releases therefore carry an estimated near-30\% label-error rate, and downstream analyses (suite-level SFR, stage distribution, ML baselines) inherit this noise. The 47-trace gold set (about 1.2\% of failures) is too small to validate the rare stages; scaling human annotation to 150--300 stratified traces is a priority. (4)~\emph{Bifurcation.} $R^2 = 0.493$ is a descriptive statistic; the bifurcation is partially mechanical and cannot be decomposed into ``mechanical'' vs.\ ``real'' components. (5)~\emph{Oracle repair.} The 64\% repair rate matches the no-op baseline, demonstrating reparability but not causal attribution. (6)~\emph{Recovery.} The 15.7\% recovery rate (\Cref{sec:res-recovery-bench}) shows that feedback signals enable partial recovery, but silent failures remain hard to correct (10.4\% vs.\ 36.7\% for loud). (7)~\emph{Statistical tests.} We use McNemar tests; multinomial mixed-effects models would be more powerful.

\textbf{Ethical considerations.} This work analyses LLM agent failures on data-science tasks. No human subjects or personal data are involved. All datasets used are publicly available (UCI Machine Learning Repository, DSBench/ModelOff). The API costs (\$80 total) were modest. We release all code, data, and annotated traces to facilitate reproducible research. Additional ethical considerations: (1)~The ChatAnywhere API endpoint may log prompts and responses; users should avoid submitting sensitive data. (2)~DSBench workbooks may contain simulated financial data; users should verify licensing terms before redistribution. (3)~The benchmark evaluates primarily closed-source models (GPT-4o, DeepSeek, Qwen), introducing a bias toward these providers; open-source models should be tested in future work. (4)~All tasks are in English and use Python, introducing language and tool bias. (5)~Low performance by certain models (e.g., DeepSeek-R1) is attributable to infrastructure failures, not fundamental model inability; results should not be used to rank models without acknowledging this caveat. (6)~Benchmark gaming risk: since synthetic tasks use fixed seeds and DSBench tasks are public, models trained on these tasks may achieve inflated scores; we recommend using the diagnostic metrics (SFR, propagation depth) rather than raw SR for model comparison.

\textbf{Generative AI usage.} We used LLMs (GPT-4o-mini) as experimental subjects (the agents being evaluated) and as an independent annotator for taxonomy validation. All LLM-generated outputs are treated as data, not as authorial content. The paper itself was written by the authors.
